\documentclass[12pt,a4paper]{article}

\usepackage{ctex}
\usepackage{geometry}
\usepackage{hyperref}
\usepackage{enumitem}
\usepackage{titlesec}

\geometry{left=2.5cm, right=2.5cm, top=2.5cm, bottom=2.5cm}

\hypersetup{
    colorlinks=true,
    linkcolor=blue,
    urlcolor=blue
}

\titleformat{\section}{\bfseries\LARGE\heiti}{\thesection}{1em}{}[\titlerule]
\titleformat{\subsection}{\bfseries\Large\heiti}{\thesubsection}{1em}{}

\begin{document}

\begin{center}
{\bfseries\Huge\heiti 弗拉基米尔·谢尔金} \\[0.5cm]
机器学习工程师 \\[0.3cm]
个人网站：\href{https://serezhinvm.github.io/serzhin-portfolio/}{\underline{serzhin-portfolio}}（点击右上角可切换至中文）\\
\end{center}

\section{教育背景}
\addcontentsline{toc}{section}{教育背景}

\subsection{圣彼得堡彼得大帝理工大学，俄罗斯}
\textit{学士 | 机电一体化与机器人技术} \hfill \textbf{2023 - 2027}

基础教育：学习机电一体化和机器人技术基础、高级编程语言、高等数学、工程学科、计算机视觉系统和人工智能。

\subsection{西安交通大学，中国}
\textit{交换学习(硕士) | 机械工程} \hfill \textbf{2026年2月 - 2027年1月}

综合技术教育。学习``编程语言理论与实现''、``软件开发''、``电气工程中的AI''、``AI与智能机器人''、``统计信号处理：基础与机器学习视角''（西安交通大学在2025年U.S. News \& World Report排名中位列机械工程专业全球第一）。

\subsection{西安工业大学，中国}
\textit{交换学习(本科) | 机械制造及其自动化} \hspace{1em} \textbf{2025年2月 - 2025年7月}

综合技术教育。核心课程：``C++程序设计'' - 93/100，``Python'' - 91/100，``机械制图'' - 100/100。

\section{专业技能}

\begin{itemize}[noitemsep]
    \item \textbf{编程语言：} Python (PyTorch, Numpy, Pandas, Matplotlib, Scikit-learn), C++, Haskell
    \item \textbf{机器人系统：} ROS, Gazebo, RViz
    \item \textbf{数学基础：} 微积分、线性代数、概率论与数理统计
    \item \textbf{其他工具：} OpenCode, 面向对象编程(基础知识)
\end{itemize}

\section{项目经验}

\subsection{MNIST手写数字识别 - CNN (PyTorch)}
用于手写数字识别的卷积神经网络。模型在MNIST数据集上训练，测试准确率达到99\%以上。包含训练过程可视化和结果分析。

技术栈：Python (PyTorch, Pandas, NumPy, Matplotlib)

\subsection{ROS机器人自主导航开发}
在ROS中实现了6个实验的组合：环境设置、Gazebo中的运动模拟、LIDAR和里程计处理、全局和局部路径规划、PID控制器调优、模糊控制实现。

技术栈：ROS, Gazebo, RViz, LIDAR, SLAM, PID, Fuzzy Logic, A*, DWA

\subsection{个人作品集网站}
机器学习工程师个人作品集网站，支持三种语言和AI聊天机器人。由AI代理OpenCode使用GitHub CLI创建。

技术栈：HTML, JavaScript, GitHub Pages

\subsection{热门数据集分析}
热门数据集分析：泰坦尼克号幸存者、美国出生率、伦敦骑自行车人数等。使用Python库进行数据处理和可视化。

技术栈：Python (Scikit-learn, Pandas, NumPy, Matplotlib, Datetime)

\subsection{水文与水资源应用任务}
使用Python、外部API、科学计算库（NumPy、SciPy）和可视化（Matplotlib、Streamlit）对水文过程进行数字建模。每个实验都包括对结果进行物理有效性验证。

技术栈：Python 3.10+, requests, Streamlit, NumPy, SciPy, Matplotlib, pandas, OpenWeatherMap API

\section{语言能力}

\begin{itemize}[noitemsep]
    \item \textbf{俄语：} 母语
    \item \textbf{英语：} C1 - 熟练（牛津分级测试）
    \item \textbf{中文：} 中级（在中国学习1.5年）
\end{itemize}

\section{证书}

\begin{itemize}[noitemsep]
    \item 成就证书 - 西安工业大学，中国
    \item 牛津分级测试 - C1
    \item 数学分析 - 在线课程（优秀）
    \item 概率论与数理统计 - 在线课程（优秀）
    \item 无人机系统软件开发 - 在线课程（144小时）
\end{itemize}

\section{联系方式}

\begin{itemize}[noitemsep]
    \item \textbf{邮箱：} svovamtr@mail.ru
    \item \textbf{Telegram：} @seriozhin
    \item \textbf{微信：} wxid\_1kay0tsy4dt722
    \item \textbf{GitHub：} \href{https://github.com/serezhinVM}{serezhinVM}
\end{itemize}

\end{document}